% 325_Hawking_FFGFT_De_ch.tex
% Dok. 325, FFGFT-Korpus, August 2026 — Kapitel-Datei
% Wird von Sources/wr_standalone_A4/325_Hawking_FFGFT_De.tex eingebunden.

% --- Dokumentlokale Makros (kollisionssicher) ---
\providecommand{\K}{}\renewcommand{\K}{\textbf{[K]}}
\providecommand{\B}{}\renewcommand{\B}{\textbf{[B]}}
\renewcommand{\S}{\textbf{[S]}}
\providecommand{\Q}{}\renewcommand{\Q}{\textbf{[Q]}}
\providecommand{\X}{}\renewcommand{\X}{\textbf{[X]}}
\providecommand{\SETZ}{}\renewcommand{\SETZ}{\textbf{[SETZUNG]}}
\providecommand{\lP}{}\renewcommand{\lP}{\ell_{\mathrm{P}}}
\providecommand{\mP}{}\renewcommand{\mP}{m_{\mathrm{P}}}
\providecommand{\TP}{}\renewcommand{\TP}{T_{\mathrm{P}}}
\renewcommand{\TH}{T_{\mathrm{H}}}
\providecommand{\rs}{}\renewcommand{\rs}{r_{\mathrm{s}}}
\providecommand{\SBH}{}\renewcommand{\SBH}{S_{\mathrm{BH}}}
\providecommand{\xig}{}\renewcommand{\xig}{\xi}
\providecommand{\Tmass}{}\renewcommand{\Tmass}{\tilde{T}}
\providecommand{\Ebit}{}\renewcommand{\Ebit}{E_{\mathrm{bit}}}
% --- Farbboxen ---
\definecolor{ffgftblue}{RGB}{30,90,160}
\definecolor{matzkegreen}{RGB}{0,120,60}
\definecolor{converge}{RGB}{160,60,0}
\definecolor{diverge}{RGB}{140,0,40}
\definecolor{warn}{RGB}{100,0,120}
\tcbset{
  base/.style={breakable,enhanced,boxrule=0.6pt,arc=3pt,
    fonttitle=\bfseries\small,coltitle=white,top=4pt,bottom=4pt,left=6pt,right=6pt},
  ffgftbox/.style={base,colback=blue!5,colframe=ffgftblue,colbacktitle=ffgftblue},
  matzbox/.style={base,colback=green!5,colframe=matzkegreen,colbacktitle=matzkegreen},
  convbox/.style={base,colback=orange!8,colframe=converge,colbacktitle=converge},
  divbox/.style={base,colback=red!6,colframe=diverge,colbacktitle=diverge},
  warnbox/.style={base,colback=violet!6,colframe=warn,colbacktitle=warn}
}

\section*{Zusammenfassung}

Dok.\ 313 (Kap.\ G) legte den thermodynamischen Rahmen der Hawking-Strahlung
in FFGFT: eine KMS-Regel temperiert Unruh, Gibbons–Hawking und Schwarze
Löcher gemeinsam \K, das Membranbild erhält sein Thermometer \S, die
Familienleiter trifft den Kosmos bei $\rs = R_H/2$ \S.
Offen blieb der \emph{mikroskopische Emissionsmechanismus}:
welche Moden entkommen, was das Quant trägt, warum Massives zurückfällt.
Dieses Dokument arbeitet ihn aus FFGFT-eigenen Bausteinen vollständig aus:

\begin{enumerate}
  \item \textbf{Temperatur}: Die KMS-Regel $T = \hbar/(k_B\tau)$ mit der
    Membran-Periode $\tau = 8\pi GM/c^3$ folgt aus $\Tmass\cdot m = 1$
    (Uhrenstauchung an der Membran, Dok.\ 313 G.2). \K
  \item \textbf{Quant}: Das Hawking-Quant ist das Bit der Skala $4\pi\rs$:
    $k_B\TH = \hbar c/(4\pi\rs)$ exakt — die systemabhängige Bit-Energie
    $\Ebit = \hbar c/L$ (Dok.\ 257) ausgewertet an der Horizontskala.
    Ein universeller Bitwert wird nirgends benötigt. \K
  \item \textbf{Selektion}: Emittierbar sind nur Trialitäts-0-Moden
    ($T_R = 0$). Confinement (Dok.\ 321) und Hawking-Selektion sind
    \emph{ein} Prinzip; die Absorption des Partnerquants ist die
    $\mathbb{Z}_3$-Projektor-Orthogonalität $P_jP_k = 0$. \B
  \item \textbf{Träger}: Der masselose Träger ist die $n_4 = 0$-Torusmode
    (räumliche Windung ohne Massenkreis-Windung). Massive Moden
    ($n_4 \neq 0$) haben endliches $\Tmass$ und fallen zurück. \K
  \item \textbf{Information}: Das emittierte Quant trägt das Sektorpaar
    $(k, -k)$ als orthogonal lesbare Quantenzahl ($\log_2 3$ Bit) —
    plus die thermische 1-nat-Flächenbuchung. Feinkörnige Entropie bleibt
    konstant (Unitarität, Dok.\ 322). Kein Informationsparadoxon. \K
\end{enumerate}

Zwei Prüfskripte, alle Assertions bestanden:
\texttt{325\_hawking\_ffgft\_mechanismus.py} (thermodynamische Kette),
\texttt{325\_hawking\_z3\_selektion.py} ($\mathbb{Z}_3$-Algebra).



% ============================================================
\section{Ausgangslage: was Dok.\ 313 offen ließ}

Dok.\ 313 Kap.\ G etablierte drei Resultate: (a) die eine KMS-Regel
$T = \hbar/(k_B\tau)$ für Unruh, Gibbons–Hawking und Schwarze Löcher \K;
(b) das Membran-Thermometer aus $\Tmass\cdot m = 1$ \S; (c) die
Familienleiter bis zum Kosmos ($\rs = R_H/2$ exakt) \S. Ebenso die
Grenze: Verdampfung ist für stellare und supermassive Löcher 56 bzw.\ 83
Größenordnungen zu langsam für den Zeitzyklus — Hawking verschärft die
grobkörnige Entropiefrage D(ii), löst sie nicht \K.

Was fehlte, war die \emph{Mikrostruktur} der Emission:
\begin{itemize}
  \item Welche Moden des $T^4/\mathbb{Z}_3$-Spektrums können den
    Horizont verlassen — und welches Prinzip selektiert sie?
  \item Was ist das emittierte Quant, und in welcher Beziehung steht
    seine Energie zur Informationseinheit des Rahmens (Dok.\ 257)?
  \item Warum fallen massive Quanten zurück — als Folge der Geometrie,
    nicht als Zusatzannahme?
  \item Wie ist die Informationserhaltung mikroskopisch realisiert?
\end{itemize}
Diese vier Fragen werden hier beantwortet — jede aus einem bereits
etablierten FFGFT-Baustein: $\Tmass\cdot m = 1$ (Dok.\ 077, 180),
systemabhängige Bit-Energie (Dok.\ 257, 302, A271–A273),
$\mathbb{Z}_3$-Trialität (Dok.\ 321), Unitarität des Spektraloperators
(Dok.\ 322).

% ============================================================
\section{Baustein 1: Die KMS-Temperatur aus \texorpdfstring{$\Tmass\cdot m = 1$}{T*m=1}}

Wo die lokale Zeitstruktur periodisch ist mit Periode $\tau$,
liest ein Beobachter Temperatur ab:
\begin{equation}
  T = \frac{\hbar}{k_B\,\tau}
  \label{eq:kms}
\end{equation}

\begin{center}
\small
\begin{tabular}{@{}lll@{}}
  \toprule
  System & Periode $\tau$ & $T = \hbar/(k_B\tau)$ \\
  \midrule
  Kosmos (kompakte Zeit) & $2\pi/H_0 = \num{2.9e18}$ s & $\num{2.66e-30}$ K \\
  Schwarzes Loch ($M_\odot$) & $8\pi GM/c^3 = \num{1.24e-4}$ s & $\num{6.17e-8}$ K \\
  Unruh ($a = g$) & $2\pi c/a = \num{1.9e8}$ s & $\num{3.98e-20}$ K \\
  \bottomrule
\end{tabular}
\end{center}

Die Membran-Periode $\tau = 8\pi GM/c^3$ ist keine Zusatzannahme,
sondern folgt aus der Zeit-Masse-Dualität: die Masse staucht die
Zeitstruktur ($\Tmass\cdot m = 1$ — je näher der Membran, desto langsamer
die Uhren), und an der Membran wird die lokale Zeitstruktur periodisch
mit genau dieser Periode (Dok.\ 313 G.2). \K

Damit ist die Temperatur nicht das Ergebnis einer Bogoliubov-Rechnung
auf gekrümmtem Hintergrund, sondern eine Ableseregel der
Uhrenlandkarte — dieselbe Regel für alle drei Systeme.

% ============================================================
\section{Baustein 2: Das Hawking-Quant als systemabhängiges Bit}

Der zentrale neue Befund dieses Dokuments. Die systemabhängige
Bit-Energie der FFGFT (Dok.\ 257) lautet $\Ebit(L) = \hbar c/L$.
Ausgewertet an der Horizontskala:

\begin{tcolorbox}[ffgftbox, title={Kernresultat \K}]
  \begin{equation}
    k_B \TH \;=\; \frac{\hbar c}{4\pi \rs}
    \;=\; \Ebit\!\left(L = 4\pi\rs\right)
  \end{equation}
  Das thermische Hawking-Quant \emph{ist} das Bit der Skala
  $L = 4\pi\rs$ (doppelter Horizontumfang).
  Numerisch exakt (Skript 1): $\hbar c/(k_B\TH) = 4\pi\rs$ auf
  Maschinengenauigkeit.
\end{tcolorbox}

Die Systemabhängigkeit der Bitwerte (Dok.\ 257, 302, A271–A273)
ist damit kein Hindernis für die Hawking-Physik, sondern
\emph{ihre Erklärung}: die Hawking-Temperatur ist die Bit-Temperatur
der Horizontskala. Ein universeller Bitwert — etwa eine feste
Bit-Masse aus einer Landauer-Auswertung bei Planck-Temperatur —
wird an keiner Stelle benötigt; jedes System trägt seine eigene
Bit-Energie, und die des Schwarzen Lochs an seinem Horizont ist
exakt $k_B\TH$.

\subsection{Flächenbilanz}

Emission eines Quants der Energie $k_B\TH$ trägt nach Clausius
genau $\dd S = E/(T\cdot k_B) = 1$ nat, also
\begin{equation}
  \dd A = -4\lP^2 \quad\text{pro nat}
\end{equation}
(Skript 1). Der Flächenverlust pro Emission folgt damit aus
Bekenstein $+$ Clausius allein — ohne weitere geometrische Annahmen. \B

% ============================================================
\section{Baustein 3: \texorpdfstring{$\mathbb{Z}_3$}{Z3}-Trialitätsselektion}

Der algebraische Kern des Mechanismus. Der Horizont trägt die
$\mathbb{Z}_3$-Orbifold-Struktur von $T^4/\mathbb{Z}_3$ (Dok.\ 321, 322).
Die drei $\mathbb{Z}_3$-Projektoren
\begin{equation}
  P_k = \frac{1}{3}\sum_{j=0}^{2}\omega^{-jk}\,\tau^j,
  \qquad \omega = e^{2\pi i/3}
\end{equation}
sind vollständig, idempotent und orthogonal (Skript 2):
\begin{equation}
  P_0 + P_1 + P_2 = \mathbb{1},\qquad
  P_j P_k = \delta_{jk}P_j.
\end{equation}

\begin{tcolorbox}[ffgftbox, title={Die Absorption \B}]
  Das einfallende Partnerquant löscht das Horizont-Voxel genau dann,
  wenn es im orthogonalen $\mathbb{Z}_3$-Sektor liegt:
  \begin{equation}
    P_j\,\bigl(P_k\,\psi\bigr) = 0 \quad (j \neq k)
  \end{equation}
  — exakt, kein Tunneln (Skript 2: $\|P_2(P_1\psi)\| < 10^{-15}$).
  Die Paarerzeugung an der Membran trennt ein $(k,-k)$-Sektorpaar;
  der einwärts fallende Partner wird durch die Projektor-Orthogonalität
  absorbiert, der auswärts laufende entkommt, sofern er die
  Trialitätsbedingung erfüllt.
\end{tcolorbox}

\subsection{Wer entkommt: \texorpdfstring{$T_R = 0$}{TR=0}}

Die Trialität einer Mode ist ihre $\mathbb{Z}_3$-Ladung $k$.
Emittierbar (asymptotisch frei) sind nur Kombinationen mit
$T_R = \sum k_i \equiv 0 \pmod 3$:

\begin{center}
\small
\begin{tabular}{@{}lcc@{}}
  \toprule
  Kombination & $T_R$ & emittierbar \\
  \midrule
  Einzelmode $k=0$ (Photon-Analogon) & 0 & ja \\
  Einzelmode $k=1$ (Quark) & 1 & nein \\
  Bilinear $(1,2)$ (Meson) & 0 & ja \\
  Bilinear $(1,1)$ & 2 & nein \\
  Trilinear $(1,1,1)$ (Baryon) & 0 & ja \\
  \bottomrule
\end{tabular}
\end{center}

\begin{tcolorbox}[ffgftbox, title={Confinement $=$ Hawking-Selektion \B}]
  Dok.\ 321 identifizierte Confinement als Trialitätsselektion $T_R = 0$.
  Dasselbe Prinzip wirkt am Horizont: nur farblose Objekte entkommen.
  Confinement und Hawking-Selektion sind \emph{ein} Prinzip
  der $\mathbb{Z}_3$-Struktur — nicht zwei.
  Die $\mathbb{Z}_3$-Ladung ist bei jeder Emission exakt erhalten
  (Skript 2).
\end{tcolorbox}

% ============================================================
\section{Baustein 4: Der masselose Träger}

Auf $T^4 = T^3_{\mathrm{Raum}}\times S^1_m$ trägt die Masse einer Mode
die Windung $n_4$ auf dem Massenkreis ($\Tmass\cdot m = 1$).
Der emittierbare Träger ist die Mode mit
\begin{equation}
  n_4 = 0:\qquad \text{räumliche Windung ohne Massenwindung.}
\end{equation}
Ein massives Quant ($n_4 \neq 0$) hat endliches $\Tmass$; in der
gestauchten Uhrenlandkarte an der Membran (Dok.\ 313 G.2) fällt es
zurück — als Folge der Geometrie, nicht als Zusatzannahme.
Bei $|n_i| \leq 1$ existieren $3^3 - 1 = 26$ masselose
Nachbarmoden gegen $54$ massive (Skript 2). \K

Das Hawking-Spektrum besteht damit aus $n_4 = 0$-Moden mit $T_R = 0$:
masselos, farblos, thermisch bei $\TH$ — die photonartigen Freiheitsgrade
des Torus.

% ============================================================
\section{Baustein 5: Informationserhaltung}

\subsection{Was das Quant trägt}

Das emittierte $T_R = 0$-Quant trägt das \emph{Sektorpaar} $(k, -k)$
als innere Quantenzahl. Die drei möglichen Paarzustände sind exakt
orthogonal (Gram-Matrix $= \mathbb{1}_3$, Skript 2) — welcher Sektor
entnommen wurde, ist am Quant lesbar:
$\log_2 3 \approx 1{,}585$ Bit Sektorinformation pro Quant,
zusätzlich zur thermischen 1-nat-Flächenbuchung.

\subsection{Feinkörnige Konstanz}

Dok.\ 313 G.5 formulierte es als \enquote{zwei Sprachen, eine Aussage}:
lokal korreliert die Strahlung mit der fraktalen Membranstruktur
($S_{\mathrm{BH}}(M_\odot) = \num{1.05e77}\,k_B$ an Korrelationen —
nichts geht verloren); global ist die feinkörnige Entropie unter
unitärer Evolution konstant (Unitarität des Spektraloperators,
Dok.\ 322 \K). Kein Informationsparadoxon: die Membran ist ein
Uhrenlandkarten-Objekt, kein kausaler Abgrund.

% ============================================================
\section{Die FFGFT-Korrektur und ihre Grenzen}

Dok.\ 313 gibt die fraktale Korrektur der Hawking-Leistung:
\begin{equation}
  P = P_{\mathrm{std}}\,\bigl(1 - \xi\ln(M/\mP)\bigr),
\end{equation}
numerisch 0{,}60\,\% (PBH $10^{12}$ kg) bis 1{,}44\,\% (SMBH $10^9\,M_\odot$)
— Skript 1. Sie ändert an den Verdampfungsbefunden nichts:
nur primordiale Löcher unterhalb
\begin{equation}
  M_* = \left(\frac{\tau_c\,\hbar\,c^4}{5120\,\pi\,G^2}\right)^{1/3}
      = \num{3.3e11}\ \mathrm{kg}
\end{equation}
schließen via Hawking innerhalb eines Zeitzyklus $\tau_c = 91{,}9$ Gyr.
Die Sonne ist 56, ein SMBH 83 Größenordnungen zu langsam
(Skript 1) — Hawking verschärft die grobkörnige Entropiefrage
D(ii), löst sie nicht (Dok.\ 313 G.4). \K

Die Familienleiter läuft dabei bruchlos bis zum Kosmos:
$\TH = T_{\mathrm{GH}}$ verlangt $M = c^3/(4GH_0) = \num{4.66e52}$ kg
mit $\rs = R_H/2$ exakt (Skript 1). \K

\begin{tcolorbox}[warnbox, title={Einordnung des Bit-Befunds}]
  $k_B\TH = \Ebit(4\pi\rs)$ zeigt: die Hawking-Physik lässt sich
  vollständig mit der systemabhängigen Bit-Energie formulieren.
  Jede Konstruktion, die stattdessen einen universellen Bitwert
  ansetzt (etwa eine feste Bit-Masse aus $k_B T\ln 2$ bei einer
  privilegierten Temperatur), fügt eine Parametrisierung hinzu,
  die für die Temperaturphysik nicht benötigt wird — konsistent
  mit A271–A273: Landauers Schranke ist eine Aussage über
  Träger-Operationen, nicht über abstrakte Information.
\end{tcolorbox}

% ============================================================
\section{Statusübersicht}

\begin{center}
\begin{tabular}{@{}lc@{}}
  \toprule
  \textbf{Element} & \textbf{Status} \\
  \midrule
  KMS-Regel temperiert Unruh/GH/BH gemeinsam & \K\ (Dok.\ 313) \\
  Membran-Periode aus $\Tmass\cdot m = 1$ & \K \\
  $k_B\TH = \Ebit(4\pi\rs)$ — Quant $=$ Horizont-Bit & \K\ (neu) \\
  Flächenquant $-4\lP^2$/nat aus Bekenstein$+$Clausius & \B\ (neu) \\
  $\mathbb{Z}_3$-Absorption $P_jP_k = 0$ & \B\ (neu) \\
  $T_R=0$-Selektion; Confinement $=$ Hawking-Selektion & \B\ (neu) \\
  Träger $= n_4$-freie Torusmode & \K\ (neu) \\
  Sektorinfo $\log_2 3$ Bit orthogonal lesbar & \B\ (neu) \\
  Informationserhalt (Membran $+$ Unitarität) & \K\ (313/322) \\
  FFGFT-Leistungskorrektur 0{,}6–1{,}4\,\% & \K\ (313) \\
  Graukörper-Faktoren, Modenspektrum-Detail & \S\ (offen) \\
  Rückreaktion auf Orbifold-Fixpunkte & \S\ (offen) \\
  \bottomrule
\end{tabular}
\end{center}

Der Hawking-Mechanismus steht damit in FFGFT auf \K/\B; offen bleiben
allein die Graukörper-Faktoren (Standard-QFT-Rechnung auf FFGFT-Hintergrund)
und die Rückreaktion der Emission auf die Orbifold-Fixpunktstruktur.

\section*{Prüfskripte}

\begin{itemize}
  \item \texttt{python/Dok325\_Skripte/325\_hawking\_ffgft\_mechanismus.py} —
    thermodynamische Kette: KMS, Membran-Thermometer, Horizont-Bit,
    Flächenbilanz, Korrektur, Verdampfung, Familienleiter,
    Informationsseite.
  \item \texttt{python/Dok325\_Skripte/325\_hawking\_z3\_selektion.py} —
    $\mathbb{Z}_3$-Algebra: Projektoren, Absorption, Trialitätsselektion,
    Ladungserhaltung, Sektorinformation, Masselosigkeit.
\end{itemize}

\section*{Literatur}

\begin{enumerate}
  \item J.\ Pascher, Dok.\ 313 \textit{Kein Anfang}, Kap.\ G (Hawking auf dem Zeitzyklus), FFGFT-Korpus, 2026.
  \item J.\ Pascher, Dok.\ 321 \textit{SU(3)$_c$ aus $\mathbb{Z}_3$-Trialität}, FFGFT-Korpus, 2026.
  \item J.\ Pascher, Dok.\ 322 \textit{Spektraltheorie und Hilbertraum-Einbettung}, FFGFT-Korpus, 2026.
  \item J.\ Pascher, Dok.\ 257, 302, A271–A273 (Bit-Energie, Landauer), FFGFT-Korpus.
  \item J.\ Pascher, Dok.\ 180 \textit{Derivation of $L_0$}, FFGFT-Korpus, 2026.
  \item S.\ W.\ Hawking, Commun.\ Math.\ Phys.\ \textbf{43}, 199, 1975.
  \item J.\ Bekenstein, Phys.\ Rev.\ D \textbf{7}, 2333, 1973.
  \item W.\ G.\ Unruh, Phys.\ Rev.\ D \textbf{14}, 870, 1976.
\end{enumerate}

\vspace{6pt}\noindent\small
\textcopyright\ 2026 Johann Pascher · \href{https://creativecommons.org/licenses/by/4.0/}{CC BY 4.0}
